This appendix gives the descriptor's rows, kernel, gap, rescaling and encoding in full.

\subsection{Block Rows}
\label{app:descriptor-rows}

\Cref{tab:app-descriptor-rows} lists every row block of the chain, on one axis and on several (state and block register: \Cref{sec:descriptor-layouts}). $\Bt_a$ and $\Dt_a$ are the banded factors on axis $a$, tensored with the identity on the other axes.

\begin{table}[!htbp]
    \centering
    \caption[Row blocks of the descriptor's chain]{\textbf{The chain's row blocks: one equation row, banded recurrences, padding and constraints.} Each row block of the descriptor's stacked residual, the state blocks it reads and its operator. $\fhat{m}$, coefficients of the $m$-th derivative; $\M_{a_j}$, multiplication by the coefficient $a_j$ (the standard form's, lifted where it lifts); $\M_r$, multiplication by the source $r$; $\Dzero(x_s)$, collapse row at the datum point $x_s$; $\Bt_a$, $\Dt_a$, the banded factors of $\G=\Bt^{-1}\Dt$ along axis $a$; $c_{\mathrm{pad}}$, padding weight; $\Brow(x^*)$, zero row at $x^*$. Constraint rows are the standard form's (\Cref{def:standard-form}) with each power of $\G$ replaced by the block that carries the derivative.}
    \label{tab:app-descriptor-rows}
    \small
    \begin{tabular}{@{}p{0.24\textwidth} p{0.24\textwidth} p{0.44\textwidth}@{}}
        \toprule
        Row block & Reads & Content \\
        \midrule
        The equation & the field and every carried derivative & $\sum_j\M_{a_j}\fhat{j}-\M_r\Dzero(x_s)\fhat{0}$ \\
        Recurrence $m\to m+1$ along axis $a$ & blocks $m$ and $m+1$ of $a$'s chain & $-\Dt_a\,\fhat{m}+\Bt_a\,\fhat{m+1}$ \\
        Padding block $b$ & block $b$ & $c_{\mathrm{pad}}\,I$, $c_{\mathrm{pad}}=1$ \\
        A constraint on $\partial_a^mf$ & block $m$ of $a$'s chain & a zero row $\Brow(x^*)$, a collapse row, a datum slice (\Cref{eq:datum-slice}) or a ratio row (\Cref{sec:systems-data}) \\
        \bottomrule
    \end{tabular}
\end{table}

A Neumann condition reads the derivative block with the zero row $\Brow(x)$, the same row that a Dirichlet condition applies to the field block, and no $\G$. A source enters the equation row as in the standard form, through the collapse row on the field block. A first-order equation or system takes the mass form of \Cref{eq:mass-form} instead of a chain, with the standard form's state and constraint rows. The system rows of several fields and of a Carleman lift, and their ratio rows, are given in \Cref{app:systems}.

\begin{proposition}[The padding decouples]
\label{prop:padding-decouples}
After ordering the columns,
\begin{equation*}
    R_{\mathrm{desc}}=\begin{pmatrix}R_{\mathrm{phys}}&0\\0&c_{\mathrm{pad}}I\end{pmatrix},
\end{equation*}
so $\Hdesc$ is block diagonal, every polynomial in $\Hdesc$ maps the physical blocks to themselves, and the padding's singular values all equal $c_{\mathrm{pad}}$.
\end{proposition}
\begin{proof}
No equation, recurrence or constraint row reads a padding block, and each padding row block is $c_{\mathrm{pad}}I$ on its own block.
\end{proof}

The reference, the gap, the filter and the initial state are therefore the physical subspace's, and $c_{\mathrm{pad}}=1$ adds one to $\alpha_R^2$ and nothing to the gap the preparation meets. Over the whole spectrum the padding's singular value can lie below the physical subspace's first excited one, but an initial state on the field block never meets it.

\subsection{Variable Coefficients and Padded Block Count}
\label{sec:app-variable-coeff}

A variable coefficient $a_j(x)$ enters the equation row through the standard form's multiplication operator $\M_{a_j}$, the lift of Paine et al.\ (\Cref{prop:lift}), on exactly the axes where the standard form lifts it: an axis with a non-constant coefficient or a source. The descriptor's equation row therefore applies the same operators as the standard form's equation block, which \Cref{thm:kernel-equivalence} needs. A multiplication by a polynomial is banded, with bandwidth the polynomial's degree, so it compiles to the same bands as $\Bt$ and $\Dt$ (\Cref{app:term-ir}). The square Galerkin form, which truncates the exact product to the first $N$ modes and never lifts, is used only as a control in the gap study of \Cref{app:descriptor-gap}.

With $K=1+\sum_ak_a$ physical blocks, the padded state has
\begin{equation}
    \dim w = 2^{\lceil\log_2K\rceil}\prod_aN_a,
    \qquad
    n_f=\lceil\log_2K\rceil ,
    \label{eq:app-padded-count}
\end{equation}
amplitudes, against the standard form's $\prod_aN_a$. On one axis $K=k+1$:

\begin{center}
\begin{tabular}{ccccc}
\toprule
Order $k$ & Blocks $k+1$ & Field qubits $n_f$ & Padded blocks $2^{n_f}$ & Padding blocks \\
\midrule
2 & 3 & 2 & 4 & 1 \\
3 & 4 & 2 & 4 & 0 \\
4 & 5 & 3 & 8 & 3 \\
7 & 8 & 3 & 8 & 0 \\
\bottomrule
\end{tabular}
\end{center}

A first-order equation takes the mass form and carries no derivative block. The heat equation has $K=4$, and Laplace's and the wave equation $K=5$, padded to eight. The padded count equals the unpadded one exactly when $K$ is a power of two and is less than twice it otherwise; the chain's state is at most $2K$ times the standard form's. This is the qubit cost that accompanies the gate counts of \Cref{sec:descriptor-cost}.

\subsection{Kernel Equivalence on Several Axes and Under Rescaling}
\label{app:proof-kernel-equivalence}

The proof of \Cref{thm:kernel-equivalence} in the body treats one axis without rescaling. Both extensions follow from one identity. Let $E$ map a field's coefficients $p$ to the state whose block of order $m$ along axis $a$ is $\G_a^mp$, with $\G_a$ the differentiation matrix on axis $a$ tensored with the identity on the others, and whose padding blocks are zero. Let $S$ be the diagonal of block weights of \Cref{sec:descriptor-rescaling}, equal to $\prod_a\zeta_a^{m_a}$ on a block of order $m_a$ along each axis and to one on the field and padding blocks.

\begin{proposition}[Substitution identity]
\label{prop:substitution-identity}
$R_{\mathrm{desc}}S\,(S^{-1}Ep)=\bigl(R_{\mathrm{std}}\,p;\,0\bigr)$, after ordering the rows: the recurrence and padding rows vanish on $Ep$, and every other row block acts on $Ep$ as the standard form's row block of the same name acts on $p$.
\end{proposition}
\begin{proof}
The recurrence of order $m$ along axis $a$ gives $-\Dt_a\G_a^mp+\Bt_a\G_a^{m+1}p=(\Bt_a\G_a-\Dt_a)\G_a^mp=0$, and the padding blocks of $Ep$ are zero. The equation row gives $\sum_j\M_{a_j}\G^jp-\M_r\Dzero(x_s)p$ term by term, each term's derivative being read from the block that holds it, which is the standard form's equation block. A constraint on $\partial_a^mf$ applies its row to $\G_a^mp$, which is the standard form's row. The factor $S$ multiplies the columns of each block by its weight and $S^{-1}$ divides the block of $Ep$ by the same weight, so they cancel.
\end{proof}

\begin{theorem}[Kernel equivalence, general form]
\label{thm:kernel-equivalence-general}
$\ker(R_{\mathrm{desc}}S)=S^{-1}E\ker R_{\mathrm{std}}$, and $E$ is injective, so the two kernels have the same dimension.
\end{theorem}
\begin{proof}
If $p\in\ker R_{\mathrm{std}}$, \Cref{prop:substitution-identity} gives $(R_{\mathrm{desc}}S)(S^{-1}Ep)=0$. Conversely, let $(R_{\mathrm{desc}}S)v=0$ and $w=Sv$. The padding rows force the padding blocks of $w$ to zero. Along each axis's chain each recurrence gives $w_{(a,m+1)}=\Bt_a^{-1}\Dt_aw_{(a,m)}=\G_aw_{(a,m)}$, since $\Bt_a$ is invertible, so $w=Ew_0$, and \Cref{prop:substitution-identity} gives $R_{\mathrm{std}}w_0=0$. $E$ is injective because its first block is the identity.
\end{proof}

\paragraph{Limits of the Mechanism.} The theorem says that an exact kernel is present in both forms or in neither, with the same dimension; a one-dimensional kernel is one-dimensional in both. It does not say that near-kernels are equally sharp. The ratio of the first excited to the lowest singular value, which the trust threshold of \Cref{sec:pipeline-reference} reads, is each form's own and differs between them. For a Carleman lift the kernel is the truncated lift's in both forms, and its uniqueness is the initial-value problem's, not a property of either form (\Cref{sec:carleman}).

\subsection{Gap Mechanism and Excluded Artefacts}
\label{app:descriptor-gap}

\begin{proposition}[The quotient on represented states]
\label{prop:descriptor-quotient}
On the represented states the Rayleigh quotient of the rescaled descriptor is
\begin{equation}
    \frac{\|R_{\mathrm{desc}}S\,(S^{-1}Ep)\|^2}{\|S^{-1}Ep\|^2}
    =\frac{\|R_{\mathrm{std}}\,p\|^2}{\sum_{m=0}^{k}\zeta^{-2m}\|\G^mp\|^2},
\end{equation}
and the $j$-th smallest singular value of $R_{\mathrm{desc}}S$ is at most the least, over $(j+1)$-dimensional subspaces $V$ of coefficient vectors, of the largest square root of this quotient on $V$.
\end{proposition}
\begin{proof}
\Cref{prop:substitution-identity} gives the numerator, and the denominator is the norm of the represented state, whose block $m$ is $\zeta^{-m}\G^mp$. By the Courant--Fischer theorem the $j$-th smallest singular value is the least, over $(j+1)$-dimensional subspaces of the state space, of the largest quotient on the subspace; restricting the subspaces to the range of $S^{-1}E$ can only raise it.
\end{proof}

The stationary values of the quotient are the generalised singular values of the pair $(R_{\mathrm{std}},S^{-1}E)$ \cite{vanloan1976gsvd}. The bound is not tight. States that break the recurrences cost little where $\Bt$'s least singular value is small, and it falls with $N$, as about $\pi/N$, with its singular vector at the ends of the interval. Measured against the bound at $n=8$, the descriptor's lowest non-zero singular value is $\res[2]{gap-control/end-zero/n8}{consistent_ratio}$ to $\res[2]{gap-control/first-order/n8}{consistent_ratio}$ times it on the controlled equations of \Cref{tab:app-gap-controls} with a single zero and on the first-order one without, $\res[2]{gap-control/constant/n8}{consistent_ratio}$ and $\res[2]{gap-control/variable/n8}{consistent_ratio}$ times it on the second-order ones without a zero, $\res[1]{gap-control/both-ends/n8}{consistent_ratio}$ times it on the equation with zeros at both ends, and $\res[2]{consistent-ratio/R5a/n4}{ratio}$ to $\res[2]{consistent-ratio/R5b/n4}{ratio}$ times it on the two-axis problems at $n=4$ per axis.

\paragraph{Concentrated States.} Let $a_k$ vanish simply at an endpoint, and take the concentrated states of \Cref{sec:descriptor-gap}. On such a state $\|a_kp^{(k)}\|=\mathcal{O}(N^{-2})\|p^{(k)}\|$ and the lower orders are of the same size, so the square root of the quotient is $\mathcal{O}(N^{-2}\zeta^k)$ against the top block. Unscaled, the lowest non-zero singular value is $\mathcal{O}(N^{-2})$ and the gap $\mathcal{O}(N^{-4})$. At an interior zero, where the resolution is $\mathcal{O}(N^{-1})$, the gap is $\mathcal{O}(N^{-2})$. At a double zero at an endpoint, $\|a_kp^{(k)}\|=\mathcal{O}(N^{-4})\|p^{(k)}\|$ and the gap is $\mathcal{O}(N^{-8})$. These are heuristics; the rates are measured. \Cref{tab:app-gap-controls} gives the unscaled gap on controlled equations whose solutions are polynomials where the coefficient vanishes, so that their kernels are exact, built with the square multiplication.

\begin{table}[!htbp]
    \centering
    \caption[Unscaled descriptor gap against zeros of the leading coefficient]{\textbf{A zero of the leading coefficient closes the unscaled descriptor's gap at the rate the Chebyshev resolution at the zero predicts.} Unscaled descriptor gap at $n=5$ and $n=9$, its factor of decrease per added qubit, and the standard form's gap, on seven equations whose leading coefficient is bounded away from zero, has a simple zero, or vanishes at both ends, each solved on the square multiplication with no lift. Exact classical computation.}
    \label{tab:app-gap-controls}
    \small
    \begin{tabular}{@{}l l l l l@{}}
        \toprule
        Equation & Leading coefficient & Descriptor, $n=5\to9$ & Per qubit & Standard \\
        \midrule
        $f''+4f'+4f$ & $1$ & $\res[3]{gap-control/constant/n5}{descriptor_gap}\to\res[3]{gap-control/constant/n9}{descriptor_gap}$ & flat & $\res[3]{gap-control/constant/n9}{standard_gap}$ \\
        $(2+x)f''+f'-f$ & $\ge1$ & $\res[3]{gap-control/variable/n5}{descriptor_gap}\to\res[3]{gap-control/variable/n9}{descriptor_gap}$ & flat & $\res[3]{gap-control/variable/n9}{standard_gap}$ \\
        $(2+x)f'+f$ & $\ge1$ & $\res[3]{gap-control/first-order/n5}{descriptor_gap}\to\res[3]{gap-control/first-order/n9}{descriptor_gap}$ & flat & $\res[3]{gap-control/first-order/n9}{standard_gap}$ \\
        $(1-x)f''-xf'+3f$ & zero at $x=1$ & $\res[2,sci]{gap-control/end-zero/n5}{descriptor_gap}\to\res[2,sci]{gap-control/end-zero/n9}{descriptor_gap}$ & $\div\res[2]{gap-control/end-zero/per-qubit}{factor}$ & $\res[3]{gap-control/end-zero/n9}{standard_gap}$ \\
        $(1-x)f'+f$ & zero at $x=1$ & $\res[2,sci]{gap-control/first-order-end-zero/n5}{descriptor_gap}\to\res[2,sci]{gap-control/first-order-end-zero/n9}{descriptor_gap}$ & $\div\res[2]{gap-control/first-order-end-zero/per-qubit}{factor}$ & $\res[3]{gap-control/first-order-end-zero/n9}{standard_gap}$ \\
        $(1-x^2)f''-2xf'+20f$ & zeros at $x=\pm1$ & $\res[2,sci]{gap-control/both-ends/n5}{descriptor_gap}\to\res[2,sci]{gap-control/both-ends/n9}{descriptor_gap}$ & $\div\res[2]{gap-control/both-ends/per-qubit}{factor}$ & $\res[3]{gap-control/both-ends/n9}{standard_gap}$ \\
        $xf''+(1-x)f'+3f$ & zero at $x=0$ & $\res[2,sci]{gap-control/interior-zero/n5}{descriptor_gap}\to\res[2,sci]{gap-control/interior-zero/n9}{descriptor_gap}$ & $\div\res[2]{gap-control/interior-zero/per-qubit}{factor}$ & $\res[3]{gap-control/interior-zero/n9}{standard_gap}$--$\res[3]{gap-control/interior-zero/n5}{standard_gap}$ \\
        \bottomrule
    \end{tabular}
\end{table}

On the equations whose leading coefficient vanishes, the first excited state holds $\res[2]{gap-control/zeros}{top_low}$ to $\res[2]{gap-control/zeros}{top_high}$ of its weight in the top-derivative block from $n=5$ to $9$, and at $n=9$ a share of $\res[2]{gap-control/zeros}{localised_low}$ to $\res[2]{gap-control/zeros}{localised_high}$ of that block's function lies within $0.02$ of the coefficient's zero. Variable coefficients alone do not close the gap; a zero does.

\paragraph{Two Axes.} The equation constrains a sum of top derivatives. Its lowest modes are equal and opposite in two top blocks (transport, Laplace's equation) or sit in the block the equation weights least (the heat equation's $\partial_x^2f$ block, weighted $1/25$), and only the slice rows penalise them: by $\mathcal{O}(1)$ against $\|f\|$, which is small against the derivative blocks. With R5b's constraints and resolution, the lowest singular value that is not a near-solution falls from $\res[2]{two-axis/transport/n2}{lowest}$ to $\res[2]{two-axis/transport/n4}{lowest}$ for transport, from $\res[2]{two-axis/heat/n2}{lowest}$ to $\res[2]{two-axis/heat/n4}{lowest}$ for the heat equation and from $\res[2]{two-axis/laplace/n2}{lowest}$ to $\res[2]{two-axis/laplace/n4}{lowest}$ for Laplace's equation as $n$ goes from $2$ to $4$ per axis, while for $f_t+f$, which has one top block, it stays at $\res[2]{two-axis/decay/n2}{lowest}$--$\res[2]{two-axis/decay/n4}{lowest}$.

\paragraph{Weights on Slice Rows.} If the slice rows are too weak against the derivative blocks, rows that penalise a high-frequency datum more are the obvious repair. Each slice row was weighted by $(m+1)^s$, with $m$ its tangential mode, a discrete trace norm of order $s$ that leaves an exact kernel unchanged, for $s$ from $0$ to $2$ in steps of $\tfrac12$, on the unscaled descriptor of Laplace's equation (R5a) and of the heat equation (R5b) at three and four qubits per axis (exact classical computation). The cost compared is $\sqrt{\|H\|/\Delta}/\gamma$, as in \Cref{app:block-rescaling}, with $\gamma$ the geometric product state's overlap with the whole state. \Cref{tab:app-trace-weights} gives the result. The least cost lies at the same weight on every panel, and there the gap rises by a larger factor than the norm, so the cost falls. A heavier weight does worse: at $s=1$ the cost is no better than unweighted on R5a and a smaller gain than the best weight on R5b. The gain does not restore the gap. From three to four qubits per axis the unweighted gap falls by a factor of order a hundred on R5a and of order ten on R5b, and the weight's factor on the gap grows far less over the same step, so the weighted gap still falls several-fold or more per qubit. Two sizes per problem fix no rate, and the weights were not combined with block rescaling.

\begin{table}[!htbp]
    \centering
    \caption[Weights on the slice rows of a two-axis descriptor]{\textbf{Weighting the slice rows raises the gap more than the norm, but the cost falls only modestly and the gap still falls with $n$.} Unscaled descriptor of Laplace's equation (R5a) and the heat equation (R5b) at $n$ qubits per axis, with each slice row weighted by $(m+1)^s$, $m$ its tangential mode. $\Delta$ is the unweighted gap ($s=0$); $s^\ast$ is the weight of least cost $\sqrt{\|H\|/\Delta}/\gamma$ among $s\in\{0,\tfrac12,\dots,2\}$; the next three columns are the gap, the norm $\|H\|$ and the cost at $s^\ast$ over their values at $s=0$; the last is the cost at $s=1$ over that at $s=0$. A cost ratio below one is a gain. Exact classical computation.}
    \label{tab:app-trace-weights}
    \small
    \begin{tabular}{@{}l c l c c c c c@{}}
        \toprule
        Problem & $n$ & $\Delta$ & $s^\ast$ & Gap & Norm & Cost & Cost at $s=1$ \\
        \midrule
        R5a & $3$ & $\res[2,sci]{build/R5a/descriptor/n3x3/unscaled}{gap}$ & $\res[1]{trace-weights/R5a/n3x3}{best_s}$ & $\res[2]{trace-weights/R5a/n3x3}{gap_ratio}$ & $\res[2]{trace-weights/R5a/n3x3}{norm_ratio}$ & $\res[2]{trace-weights/R5a/n3x3}{cost_ratio}$ & $\res[2]{trace-weights/R5a/n3x3}{cost_ratio_s1}$ \\
        R5a & $4$ & $\res[2,sci]{build/R5a/descriptor/n4x4/unscaled}{gap}$ & $\res[1]{trace-weights/R5a/n4x4}{best_s}$ & $\res[2]{trace-weights/R5a/n4x4}{gap_ratio}$ & $\res[2]{trace-weights/R5a/n4x4}{norm_ratio}$ & $\res[2]{trace-weights/R5a/n4x4}{cost_ratio}$ & $\res[2]{trace-weights/R5a/n4x4}{cost_ratio_s1}$ \\
        R5b & $3$ & $\res[2,sci]{build/R5b/descriptor/n3x3/unscaled}{gap}$ & $\res[1]{trace-weights/R5b/n3x3}{best_s}$ & $\res[2]{trace-weights/R5b/n3x3}{gap_ratio}$ & $\res[2]{trace-weights/R5b/n3x3}{norm_ratio}$ & $\res[2]{trace-weights/R5b/n3x3}{cost_ratio}$ & $\res[2]{trace-weights/R5b/n3x3}{cost_ratio_s1}$ \\
        R5b & $4$ & $\res[2,sci]{build/R5b/descriptor/n4x4/unscaled}{gap}$ & $\res[1]{trace-weights/R5b/n4x4}{best_s}$ & $\res[2]{trace-weights/R5b/n4x4}{gap_ratio}$ & $\res[2]{trace-weights/R5b/n4x4}{norm_ratio}$ & $\res[2]{trace-weights/R5b/n4x4}{cost_ratio}$ & $\res[2]{trace-weights/R5b/n4x4}{cost_ratio_s1}$ \\
        \bottomrule
    \end{tabular}
\end{table}

\paragraph{Field-Block Weight.} Let $P_0$ project onto the field block and let the initial state be supported there. Then $\gamma=|\braket{\psi_g|\psi_{\mathrm{init}}}|=|\braket{P_0\psi_g|\psi_{\mathrm{init}}}|\le\|P_0\psi_g\|=w_f^{1/2}$ by the Cauchy--Schwarz inequality, with $w_f=\|P_0\psi_g\|^2$ the field block's weight that the build records hold. For an exact kernel, $\psi_g\propto S^{-1}Eb$ with $b$ the solution's coefficients, so $w_f=\|b\|^2/\sum_m\zeta^{-2m}\|\G^mb\|^2$. Amplitude amplification then costs a factor $1/\gamma\ge w_f^{-1/2}$ (\Cref{sec:pipeline-filter}).

\paragraph{Limits of the Weight.} Each candidate artefact was tested:

\begin{center}
\small
\begin{tabular}{@{}p{0.2\textwidth} p{0.38\textwidth} p{0.34\textwidth}@{}}
\toprule
Candidate & Test & Outcome \\
\midrule
The lift & R3a with the square multiplication & collapses faster: $\res[2,sci]{lift-artefact/R3a/n8}{square_gap}$ at $n=8$, against $\res[2,sci]{lift-artefact/R3a/n8}{lift_gap}$ lifted \\
The padding & lowest singular value against $c_{\mathrm{pad}}$; the excited state's padding weight & at most $\res[2]{padding/unscaled}{sigma1_max}$, below $c_{\mathrm{pad}}=1$, on every unscaled padded record; no padding weight \\
The constraints & controlled equations with none & they collapse too \\
Precision & lowest singular value against the decomposition's floor $\varepsilon\sigma_{\max}$ & six orders of magnitude above it at the smallest \\
The basis & a descriptor carrying $f^{(j)}$ in the ultraspherical basis $C^{(j)}$ & the same rates where $a_k$ vanishes \\
The reference & dense decomposition against the iterative reference, on R2a at $n=8$ to $10$ & agree to $\res[2]{iterative/R2a/descriptor/range}{high}\,\varepsilon\sigma_{\max}$ \\
\bottomrule
\end{tabular}
\end{center}

The ultraspherical descriptor keeps a flat gap on regular problems with a norm $\res[2]{ultraspherical/constant/n9}{norm_ratio}$ times smaller at $n=9$, and on Legendre's equation its gap is $\res[3]{ultraspherical/both-ends/n9}{gap_ratio}$ times the Chebyshev descriptor's there, the share of the recurrences, but it still collapses.

\subsection{Gap on One Axis}
\label{app:flat-gap}

This section proves \Cref{thm:flat-gap}. Write $\|g\|_T^2=\int_{-1}^1g^2(1-x^2)^{-1/2}\,dx$ and $\|g\|_U^2=\int_{-1}^1g^2(1-x^2)^{1/2}\,dx$, the norms under which the Chebyshev polynomials of the first and the second kind are orthogonal. By \Cref{app:eta-e}, a polynomial $g$ of degree below $N$ with coefficients $\hat g$ in the normalised basis has $\|\hat g\|=\sqrt{N/\pi}\,\|g\|_T$.

\begin{lemma}[The factors in the second kind]
\label{lem:second-kind}
Let $\psi_0=U_0/(2\sqrt N)$ and $\psi_m=U_m/\sqrt{2N}$ for $1\le m<N$. For a polynomial $g$ of degree below $N$, $\Bt\hat g$ and $\Dt\hat g$ are the coefficients of $g$ and of $g'$ in the basis $(\psi_m)$, and a polynomial $h=\sum_mc_m\psi_m$ has $\|c\|\ge2\sqrt{N/\pi}\,\|h\|_U$.
\end{lemma}
\begin{proof}
With $\varphi_0=T_0/\sqrt N$ and $\varphi_\ell=T_\ell\sqrt{2/N}$, the identities $T_0=U_0$, $T_1=U_1/2$ and $T_\ell=(U_\ell-U_{\ell-2})/2$ for $\ell\ge2$ give $\varphi_0=2\psi_0$, $\varphi_1=\psi_1$, $\varphi_2=\psi_2-\sqrt2\,\psi_0$ and $\varphi_\ell=\psi_\ell-\psi_{\ell-2}$ for $\ell\ge3$, the columns of $\Bt$ (\Cref{eq:app-normalised-entries}). The identity $T_\ell'=\ell U_{\ell-1}$ (\Cref{eq:app-ot-diff}) gives $\varphi_1'=2\sqrt2\,\psi_0$ and $\varphi_\ell'=2\ell\,\psi_{\ell-1}$ for $\ell\ge2$, the columns of $\Dt$. With $x=\cos\theta$, $U_m=\sin(m+1)\theta/\sin\theta$, so $\int_{-1}^1U_mU_{m'}(1-x^2)^{1/2}dx=\int_0^\pi\sin(m+1)\theta\,\sin(m'+1)\theta\,d\theta=\tfrac\pi2\delta_{mm'}$: the $\psi_m$ are orthogonal under $\|\cdot\|_U$ with $\|\psi_m\|_U^2\le\pi/(4N)$, which gives the bound.
\end{proof}

A recurrence row therefore measures the defect $F_{j+1}-F_j'$ between the functions of consecutive blocks in $\|\cdot\|_U$, a norm weaker than the state's $\|\cdot\|_T$ near the ends of the interval; this is why $\Bt$'s least singular value falls as $N$ grows. Integration recovers what the weaker norm loses. For integrable $g$ and $-1\le y\le x\le1$ the Cauchy--Schwarz inequality gives
\begin{equation}
    \Bigl|\int_y^xg\Bigr|\le\|g\|_T\,|x-y|^{1/2},
    \qquad
    \Bigl|\int_y^xg\Bigr|\le\sqrt\pi\,\|g\|_U ,
    \label{eq:app-integral-bounds}
\end{equation}
since $\int_y^x(1-t^2)^{1/2}dt\le x-y$ and $\int_{-1}^1(1-t^2)^{-1/2}dt=\pi$.

\begin{proof}[Proof of \Cref{thm:flat-gap}]
\emph{Reductions.} The padding decouples (\Cref{prop:padding-decouples}), so only the physical columns count. The rescaled residual is $R_{\mathrm{desc}}S$ with $S$ diagonal, invertible and independent of $N$ at a fixed $\zeta$; writing $\|R_{\mathrm{desc}}Sv\|/\|v\|=(\|R_{\mathrm{desc}}u\|/\|u\|)(\|u\|/\|v\|)$ with $u=Sv$, the Courant--Fischer theorem places each singular value of $R_{\mathrm{desc}}S$ between $\min_m\zeta^m$ and $\max_m\zeta^m$ times the unscaled one, so it suffices to treat $R=R_{\mathrm{desc}}$ unscaled. Divide every row by $\sqrt{N/\pi}$. A state $w$ whose blocks are the coefficients of polynomials $F_0,\dots,F_k$ of degree below $N$ then has $\|w\|^2=\sum_j\|F_j\|_T^2$. By \Cref{lem:second-kind}, by \Cref{prop:lift}, which makes the lifted equation row exact, and by $\Brow(x)b=\sqrt N f_N(x)\ket0$ (\Cref{app:eta-e}),
\begin{equation}
    \|Rw\|^2\ge4\sum_{j<k}\|F_{j+1}-F_j'\|_U^2+\|E\|_T^2+\pi\sum_i|\ell_i(w)|^2 ,
    \label{eq:app-rows-below}
\end{equation}
with $E=\sum_ja_jF_j-r_NF_0(x_s)/f(x_s)$ the equation row's function, $r_N$ the source as represented, and $\ell_i(w)$ the $i$-th constraint row's combination of the values $F_j(x^*)$, $j<k$, and $F_0(x_s)$.

\emph{The lower bound.} Let $V_N$ be the states with $\langle F_0,f\rangle_T=0$, a subspace of codimension one, so that $\sigma_2\ge\min_{w\in V_N}\|Rw\|/\|w\|$ by the Courant--Fischer theorem. Suppose no $c>0$ bounds this minimum below for every large $N$. Then there are $N\to\infty$ and $w^N\in V_N$ with $\|w^N\|=1$ and $\|Rw^N\|\to0$; write $F_j^N$ for their blocks. For $j<k$,
\begin{equation*}
    F_j^N(x)=F_j^N(0)+\int_0^xF_{j+1}^N-\int_0^x\bigl(F_{j+1}^N-(F_j^N)'\bigr).
\end{equation*}
By \Cref{eq:app-integral-bounds} the middle term is bounded by $1$ and equicontinuous, since $\|F_{j+1}^N\|_T\le1$, and the last tends to zero uniformly, since its integrand's $\|\cdot\|_U$ does by \Cref{eq:app-rows-below}. The constants are bounded as well, $\sqrt\pi\,|F_j^N(0)|\le\|F_j^N\|_T+\sqrt\pi\sup|F_j^N-F_j^N(0)|$. By the Arzelà--Ascoli theorem a subsequence has $F_j^N\to F_j$ uniformly for every $j<k$. On it the equation row gives
\begin{equation*}
    a_kF_k^N=E^N-\sum_{j<k}a_jF_j^N+r_N\frac{F_0^N(x_s)}{f(x_s)}\;\longrightarrow\;-\sum_{j<k}a_jF_j+r\,\frac{F_0(x_s)}{f(x_s)}
\end{equation*}
in $\|\cdot\|_T$, because $\|E^N\|_T\to0$ and $r_N\to r$ uniformly for an analytic source \cite[Ch.~8]{trefethen2013approximation}. Since $|a_k|$ is bounded below, $F_k^N$ converges in $\|\cdot\|_T$ to the continuous function $F_k$ that the limit defines. Letting $N\to\infty$ in the integral identity, with the first bound of \Cref{eq:app-integral-bounds} applied to $F_{j+1}^N-F_{j+1}$, gives $F_j(x)=F_j(0)+\int_0^xF_{j+1}$, so $F_0\in C^k[-1,1]$ with $F_0^{(j)}=F_j$. The limit of the display says that the equation row vanishes on $F_0$, and $\ell_i(w^N)\to0$ with uniform convergence says that every constraint row does. By hypothesis $F_0=tf$, and $\langle F_0,f\rangle_T=\lim\langle F_0^N,f\rangle_T=0$ forces $t=0$, so every $F_j$ vanishes. The convergence is in norm, however, so $\sum_j\|F_j\|_T^2=\lim\|w^N\|^2=1$, a contradiction.

\emph{The upper bound and the ground value.} Let $b$ be the coefficients of the Chebyshev truncation $f_N$ and $\hat q$ those of a fixed polynomial $q$ that is not a multiple of $f$. For $p=\alpha b+\beta\hat q$, \Cref{prop:substitution-identity} gives $\|REp\|=\|R_{\mathrm{std}}p\|$, and $\|Ep\|\ge\|p\|$. For $N>\deg q$ the rows of $R_{\mathrm{std}}$ act on $q$ through its derivatives and values, which do not depend on $N$, and through $r_N$; they act on $f_N$ through $f_N^{(j)}$, which converge to $f^{(j)}$ uniformly for $j\le k$, since the coefficients of an analytic function decay geometrically \cite[Ch.~8]{trefethen2013approximation} and $\|T_n^{(j)}\|_\infty\le n^{2j}$ by Markov's inequality \cite{borweinerdelyi1995}. So $\|R_{\mathrm{std}}p\|/\sqrt{N/\pi}$ is bounded for $|\alpha|^2+|\beta|^2=1$, as in the normalisation of the lower bound, while $\|p\|/\sqrt{N/\pi}$, which tends to $\|\alpha f+\beta q\|_T$, stays away from zero, and the Courant--Fischer theorem gives $\sigma_2\le C$. Because $f$ satisfies its equation and constraints exactly, the same convergence gives $\sigma_1\le\|R_{\mathrm{std}}b\|/\|b\|\to0$. Hence $\Delta=\sigma_2^2-\sigma_1^2$ lies between $c^2/2$ and $C^2$ for large $N$.

\emph{The mass form.} Its rows are the standard form's equation rows multiplied by $I_C\otimes\Bt$, so by \Cref{lem:second-kind} they are the coefficients of $z'-\Acoup z$ in the basis $(\psi_m)$, and \Cref{eq:app-rows-below} holds with $4\|z'-\Acoup z\|_U^2$ in place of the recurrence and equation terms. With $\delta^N=(z^N)'-\Acoup z^N$, the identity $z^N(x)=z^N(0)+\int_0^x\Acoup z^N+\int_0^x\delta^N$ and the same argument give the lower bound. Its state is the standard form's, and $\|R_{\mathrm{mass}}p\|\le(2+\sqrt2)\|R_{\mathrm{std}}p\|$ since $\|\Bt\|_2\le2+\sqrt2$, so the upper bound and the ground value follow as above.
\end{proof}

The argument gives no value for the bound, because it closes by compactness. It needs the leading coefficient's lower bound at one step only, to recover the top block from the equation row; where $a_k$ vanishes that step fails, which is the mechanism of \Cref{app:descriptor-gap}. On two axes the equation row controls only a sum of top derivatives, and no single block can be recovered from it.

\subsection{Block Rescaling}
\label{app:block-rescaling}

The rule of \Cref{sec:descriptor-rescaling} returns each axis's weight from the equation's coefficients and the resolution alone:

\begin{center}
\small
\begin{tabular}{@{}p{0.42\textwidth} p{0.5\textwidth}@{}}
\toprule
Problem & $\zeta$ \\
\midrule
One axis, regular ($a_k$ with no zero on $[-1,1]$) & $2(\bar a_j/\bar a_k)^{1/(k-j)}$, $j$ the lowest order present, $\bar a$ the mean of $|a|$ \\
One axis, singular ($a_k$ with a zero on $[-1,1]$) & $N/2$ \\
Several axes & $N_a$, times $\sqrt{\max_b\bar c_b/\bar c_a}$ on an axis of top order two, $\bar c_a$ the mean magnitude of its top coefficient \\
\bottomrule
\end{tabular}
\end{center}

Whether $a_k$ vanishes is decided from its roots, to a tolerance loose enough to find a double root. On R2a the rule gives $\zeta=2\sqrt{4/1}=4$, on R2c $\zeta=2\sqrt{400}=40$, and on R3a, whose leading coefficient vanishes at both ends, $\zeta=N/2$.

\paragraph{Justification of the Constants.} For a regular problem (\Cref{def:regular}) with constant coefficients, $(\bar a_0/\bar a_k)^{1/k}$ is the geometric mean of the characteristic roots' magnitudes, the product of the roots being $\pm a_0/a_k$; a solution $e^{\lambda x}$ has $\|f^{(m)}\|\approx|\lambda|^m\|f\|$, so dividing block $m$ by $\zeta^m$ balances the blocks at the equation's own scale. The factor two is where the cost's minimum over $\zeta$ sits on the regular benchmark problems. On a singular problem the concentrated states of \Cref{app:descriptor-gap} have quotient $\mathcal{O}(N^{-4}\zeta^{2k})$, so for second-order problems $\zeta\propto N$ makes it of order one; half the resolution is the constant the sweep found. At an interior simple zero the same argument with Bernstein's inequality in place of Markov's, states concentrated within $\mathcal{O}(N^{-1})$ and derivatives growing by $\mathcal{O}(N)$ each, gives a quotient of order $N^{-2}\zeta^{2k}$ and so $\zeta\propto N^{1/k}$, which is $N^{1/2}$ at second order; the rule, which tests only whether $a_k$ vanishes on $[-1,1]$, would give such a problem $N/2$, a weight larger than this balance needs and so a larger norm than necessary. No benchmark problem has an interior zero and this case was not run. On several axes the resolution is used per axis, as for a singular problem, and an axis of top order two is weighted so that its top-order term $\bar c_a\zeta_a^2$ balances the others: the wave equation $f_{tt}=c^2f_{xx}$, with wave speed $c$, takes $(cN,N)$, which costs at most $\res[3]{zeta-wave/R5c/n4x4}{rule_ratio}$ times the best pair on a grid at three and four qubits per axis, where $(N/2,N/2)$ costs $\res[2,sci]{zeta-wave/R5c/n3x3}{resolution_ratio}$ and $\res[2,sci]{zeta-wave/R5c/n4x4}{resolution_ratio}$ times it.

\paragraph{Calibration and Check.} The cost measured is $\sqrt{\|H\|/\Delta}/\gamma$, the walk's queries at an ideal encoding times amplitude amplification, up to the logarithms both forms share. The two constants were selected by calibration: $\zeta$ was swept over powers of two on the benchmark problems, each constant set where the cost is least, and both were then held fixed and never refitted. The same problems then checked the resulting weights, so the check is in-sample. \Cref{tab:app-zeta-window} gives the window of $\zeta$ over which the cost stays within twice its least, and \Cref{tab:app-zeta-rule} the rule's cost against the grid's best.

\begin{table}[!htbp]
    \centering
    \caption[Window of block weights near the least cost]{\textbf{The block weight need only be right to about a factor of two.} Range of the block weight $\zeta$ over which the cost $\sqrt{\|H\|/\Delta}/\gamma$ stays within twice its least, at $n=7$ on one axis and $n=4$ per axis on two (one weight shown per two-axis problem). $\|H\|$, Hamiltonian norm; $\Delta$, spectral gap; $\gamma$, the initial state's overlap; the cost counts the walk's queries at an ideal encoding times amplitude amplification. Exact classical computation.}
    \label{tab:app-zeta-window}
    \small
    \begin{tabular}{@{}lccccccc@{}}
        \toprule
        & R2a & R2c & R3a & R4a & R4d & R5a & R5b \\
        \midrule
        Window of $\zeta$ & $\res[2]{zeta-window/R2a/n7}{low}$--$\res[2]{zeta-window/R2a/n7}{high}$ & $\res[2]{zeta-window/R2c/n7}{low}$--$\res[2]{zeta-window/R2c/n7}{high}$ & $\res[2]{zeta-window/R3a/n7}{low}$--$\res[2]{zeta-window/R3a/n7}{high}$ & $\res[2]{zeta-window/R4a/n7}{low}$--$\res[2]{zeta-window/R4a/n7}{high}$ & $\res[2]{zeta-window/R4d/n7}{low}$--$\res[2]{zeta-window/R4d/n7}{high}$ & $\res[2]{zeta-window/R5a/n4x4}{low}$--$\res[2]{zeta-window/R5a/n4x4}{high}$ & $\res[2]{zeta-window/R5b/n4x4}{low}$--$\res[2]{zeta-window/R5b/n4x4}{high}$ \\
        \bottomrule
    \end{tabular}
\end{table}

\begin{table}[!htbp]
    \centering
    \caption[The block weight against a grid search]{\textbf{The block weight of \Cref{eq:zeta-rule} lands near the best on a grid.} Cost $\sqrt{\|H\|/\Delta}/\gamma$ at the rule's $\zeta$, over the least cost on a grid of powers of two; $\Delta$, gap on the physical subspace; $\gamma$, overlap of the geometric initial state on the field block. A ratio below one is the grid's spacing. Exact classical computation, in-sample: the rule's two constants were calibrated on these problems.}
    \label{tab:app-zeta-rule}
    \small
    \begin{tabular}{@{}l c l@{}}
        \toprule
        Problems & $n$ & Rule's cost / grid's best \\
        \midrule
        R1, R2a--c, R4b--d (regular) & 5, 7 & $\res[2]{zeta-rule/regular/n5}{low}$--$\res[3]{zeta-rule/regular/n5}{high}$ at $n=5$, the highest R2c's; $\res[2]{zeta-rule/regular/n7}{low}$--$\res[3]{zeta-rule/regular/n7}{high}$ at $n=7$ \\
        R3a, R3b, R4a (a simple zero) & 5, 7 & $\res[3]{zeta-rule/simple-zero/n5}{low}$--$\res[3]{zeta-rule/simple-zero/n7}{high}$ \\
        R3c, R3d (solved for $v=f/\sqrt{1-x^2}$) & 5, 7 & $\res[3]{zeta-rule/substituted/n5}{low}$--$\res[3]{zeta-rule/substituted/n5}{high}$ at $n=5$; $\res[3]{zeta-rule/substituted/n7}{low}$ and $\res[3]{zeta-rule/substituted/n7}{high}$ at $n=7$ \\
        R5a, R5b, R5c & 3, 4 per axis & $\res[3]{zeta-rule/two-axis/n3x3}{low}$--$\res[3]{zeta-rule/two-axis/n4x4}{high}$ \\
        \bottomrule
    \end{tabular}
\end{table}

\paragraph{Dip in a Coefficient.} $(\delta+x^2)f''+f=0$ has no zero but dips to $\delta$. The rule is within $\res[2]{zeta-dip/delta=0.01/n7}{ratio}$ times the grid's best for $\delta\ge0.01$ at $n=5$ and $7$, and $\res[2]{zeta-dip/delta=0.001/n5}{ratio}$ and $\res[2]{zeta-dip/delta=0.001/n7}{ratio}$ times at $\delta=10^{-3}$, where the resolution sees the dip and it behaves like a zero, so that the regular weight is the wrong one. The two classes are asymptotic, and the weight jumps between them while the coefficient changes continuously. None of the benchmark problems dips so.

\paragraph{Double Zero.} Wu et al.'s R3c and R3d multiply the associated Legendre operator by $1-x^2$, so their leading coefficient $(1-x^2)^2$ vanishes twice at each end. Unscaled, the descriptor's gap there falls as $N^{-8}$. Balancing the concentrated states needs $\zeta\propto N^2$, by which $\|\Hdesc\|$ has grown as $\zeta^4=N^8$. On a grid of $\zeta$ the best cost of the descriptor grows as about $N^{\res[2]{double-zero/R3d/faithful}{best_exponent}}$ on R3d and $N^{\res[2]{double-zero/R3c/faithful}{best_exponent}}$ on R3c over $n=5$ to $8$, against the standard form's $N^{\res[2]{double-zero/R3d/faithful}{standard_exponent}}$ and $N^{\res[2]{double-zero/R3c/faithful}{standard_exponent}}$, and at $n=8$ it is $\res[2]{double-zero/R3d/faithful}{best_over_standard}$ and $\res[2]{double-zero/R3c/faithful}{best_over_standard}$ times the standard form's; no block weight brings the two level. Solved for $v=f/\sqrt{1-x^2}$, whose leading coefficient has simple zeros, the two forms' costs grow alike, as $N^{\res[2]{double-zero/R3d/corrected}{rule_exponent}}$ and $N^{\res[2]{double-zero/R3c/corrected}{rule_exponent}}$ at the rule's weight against the standard form's $N^{\res[2]{double-zero/R3c/corrected}{standard_exponent}}$. The ground state stays accurate in both forms at the double zero, its error set by the solution's branch point; it is the cost that the double zero raises.

\paragraph{Effect on the Circuit.} The weights multiply the bands before the block register's patterns are expanded in Pauli words (\Cref{alg:ir-to-lcu}, line 3), so they change the coefficients $c_t$ and $\alpha_A=\sum_t|c_t|$ and, where a merged pattern's entries were equal across its blocks, the length of its expansion: the recurrences' $\operatorname{diag}(1,1,0,0)$, which is $\tfrac12(I+Z)$ on the block register's leading qubit, becomes $\operatorname{diag}(1,\zeta,0,0)$, which needs all four words of the two qubits. No register changes width, since a pattern on $n_f$ qubits has at most $4^{n_f}$ words whatever its entries, so the rescaling adds no ancilla; the terms, and the preparations and block-register operations that apply them, grow by at most that factor, which the equation fixes and $n$ does not. A general diagonal equilibration, a different weight for each amplitude, would need its own diagonal block encoding, with $\Theta(\dim w)$ gates, and is excluded from every circuit.

\subsection{Term IR and Compilation}
\label{app:term-ir}

The system rows of the descriptor (the equation, the recurrences and the padding; the rank-one constraint rows are encoded apart) form one square operator on the block register and the axes' registers. Each row block sits at a block index of its own: a recurrence at the index of its lower block, the equation at the lowest index left, and each padding row at its own. The term IR writes this operator as
\begin{equation}
    A_{\mathrm{sys}}=\sum_tc_t\,P_t\otimes\bigotimes_aF_{t,a},
    \label{eq:app-term-ir}
\end{equation}
with the block word $P_t$ a Pauli word $X^xZ^z$ on the block register (or, for a system's coupling, a matching, \Cref{sec:ext-circuits}), and each $F_{t,a}$ a \emph{band}: a diagonal $D$ followed by the truncated shift that takes column $q$ to row $q-\ell$ and discards it where that row does not exist, with $D$ the identity, one $Z_j$ or a reflection $Q_m$.

\begin{lemma}[An affine band with exceptions]
\label{lem:affine-band}
A diagonal whose entries are $v_q=a+bq$ except at a few columns $m$, where $v_m=a+bm+d_m$, is
\begin{equation*}
    \operatorname{diag}(v)=\Bigl(a+b\,\frac{N-1}{2}+\sum_m\frac{d_m}{2}\Bigr)I-b\sum_{j=0}^{n-1}2^{j-1}Z_j-\sum_m\frac{d_m}{2}\,Q_m :
\end{equation*}
$n+1$ terms and one reflection for each exception.
\end{lemma}
\begin{proof}
$\Lambda=\tfrac{N-1}{2}I-\sum_j2^{j-1}Z_j$ (\Cref{eq:lambda-lcu}), and $\ket{m}\bra{m}=(I-Q_m)/2$.
\end{proof}

\begin{lemma}[The affine part is tight]
\label{lem:affine-tight}
If $a$ and $a+b(N-1)$ have one sign, the combination of \Cref{lem:affine-band} without exceptions has $\alpha=\max(|a|,|a+b(N-1)|)=\|aI+b\Lambda\|_2$.
\end{lemma}
\begin{proof}
$\alpha=|a+b\tfrac{N-1}{2}|+|b|\tfrac{N-1}{2}$, the mean of the end values plus half their range, which is the larger end value in magnitude when the two share a sign.
\end{proof}

Every band of the descriptor has this form. $\Dt$'s diagonal is $2\Lambda$ except at column~1, where the normalisation of $T_0$ makes it $2\sqrt2$; $\Bt$'s and a constant multiplication's are constant except at their ends; a variable coefficient's multiplication bands and their truncation at the register's edge give the other exceptions. The columns a shift discards are free, and are fitted to keep the combination shortest.

\paragraph{Block Register.} Terms whose axis factors agree are merged across the row and column blocks they join, so that the block register carries their pattern as one small matrix $F$, encoded once; the recurrences share $\Dt$ in this way. $F$ is expanded in Pauli words, $F=\sum_{x,z}f_{xz}X^xZ^z$ with $f_{xz}=2^{-n_f}\sum_b(-1)^{z\cdot b}F_{b\oplus x,\,b}$. Where the merged pattern is a single matrix unit on one block, its $2^{n_f}$ Pauli words are replaced by one term with a flag qubit post-selected to keep that block, at the same $\alpha$. A large irregular coupling, a Carleman generator above all, is split into matchings instead (\Cref{sec:ext-circuits}).

\begin{proposition}[Compilation]
\label{prop:term-ir-compilation}
Let PREPARE$_R$ load $\operatorname{sgn}(c_t)\sqrt{|c_t|/\alpha_A}$ and PREPARE$_L$ load $\sqrt{|c_t|/\alpha_A}$ at each term's joint index of the factored selectors, with $\alpha_A=\sum_t|c_t|$. Then the circuit that applies PREPARE$_R$, then SELECT, then PREPARE$_L^\dagger$, post-selected on the selectors, the overflow qubits and the flag, block-encodes $A_{\mathrm{sys}}$ at $\alpha_A$.
\end{proposition}
\begin{proof}
On the joint index of term $t$, the factored SELECT applies the tensor product of the operations its registers hold: the term's block-register word and, on each axis, its diagonal and then its shift, which is the band. A joint index that no term uses has zero amplitude in both preparations. The post-selected block is $\sum_t\sqrt{|c_t|/\alpha_A}\cdot\operatorname{sgn}(c_t)\sqrt{|c_t|/\alpha_A}\,U_t=\sum_tc_tU_t/\alpha_A$. With one preparation and its own inverse the signs would square away, and the block would be $\sum_t|c_t|U_t/\alpha_A$.
\end{proof}

\paragraph{Truncated Shift.} On the output register widened by one overflow qubit, initially zero, add $-\ell$ modulo twice the register's size. A row $q-\ell$ inside $[0,N)$ lands in the lower half, and one outside it lands in the upper half, which the projection of the overflow qubit on $\ket0$ discards. The addition is an increment of the qubits from $j$ upwards for each set bit $j$ of $|\ell|$, each a ladder of \gls{mcx} gates, and the overflow qubit is the shift's only ancilla.

\begin{algorithm}[htbp]
    \caption{Compiling the rows of $\Adesc$ into two preparations and one SELECT.}
    \label{alg:ir-to-lcu}
    \KwIn{the blocks of $\Adesc$ as bands, the block weights $\zeta$}
    \KwOut{PREPARE$_R$, SELECT and PREPARE$_L^\dagger$; the subnormalisation $\alpha_A$}
    Write every block as a sum of bands: a diagonal and a truncated shift\;
    Fit each diagonal as affine in the column index plus a few exceptions (\Cref{lem:affine-band})\;
    Weight each band by $\zeta^{m}$ for its column block $m$\;
    Merge bands with equal axis factors across blocks; expand each resulting block-register pattern in Pauli words, giving the terms $c_tU_t$\;
    $\alpha_A \leftarrow \sum_{t} |c_t|$; on the joint selector index of term $t$, PREPARE$_R$ loads $\operatorname{sgn}(c_t)\sqrt{|c_t|/\alpha_A}$ and PREPARE$_L$ loads $\sqrt{|c_t|/\alpha_A}$\;
    SELECT: controlled on each selector register, apply the block word, then on each axis the diagonal and the shift\;
    \Return{PREPARE$_R$, SELECT, PREPARE$_L^\dagger$, $\alpha_A$}\;
\end{algorithm}

\paragraph{Rank-One Terms.} A source $\ket{\hat r}\bra{\tau(x_s)}/f(x_s)$ is compiled apart and summed with the bands by a further combination: a matrix unit on the block register, the feature map reading $\tau(x_s)$, and a rotation tree loading $\hat r$, with two ancilla and no arithmetic. The constraint rows are encoded the same way, each as its own rows, and stacked into the reflection of \Cref{sec:pipeline-walk}.

\paragraph{IR Check.} The terms are built from the problem, the operator library, the residual's layout, its lifts and its block weights, never from the assembled matrices. Summing them back and comparing the sum with the independently assembled system rows is the IR mode of operator verification, possible at any $n$ (\Cref{tab:verification-modes}); the compiled circuit is then compared with the terms exactly, or by probe states where it is too large.

\subsection{Proof of the Cost Proposition and Norms}
\label{app:proof-descriptor-cost}

\paragraph{Ancilla.} The registers are those of \Cref{app:term-ir}: one overflow qubit for each axis that shifts, at most one flag, and the factored selectors. On axis $a$ the diagonal selector indexes the identity, the $n_a$ gates $Z_j$ and the distinct reflections $r_a$, so it has $\lceil\log_2(1+n_a+r_a)\rceil$ qubits. The reflections are bounded in number by the problem's bands, and the block-register words and the shifts depend on the equation's structure, not on $n$, so every other selector has a width fixed by the problem, and the total is $\mathcal{O}(\log n)$. For R2 the reflections are three ($\Bt$'s diagonal and second superdiagonal at their first entries and $\Dt$'s first entry), so the diagonal selector has $\lceil\log_2(n+4)\rceil$ qubits; the shifts are $0$, $1$ and $2$, two qubits; the block-register selector has four; with the overflow qubit and the flag the count is $\lceil\log_2(n+4)\rceil+8$. It is verified on the compiled circuit of R2's system rows at $n=2$, $4$, $5$ and $12$.

\paragraph{Gates.} SELECT applies each distinct operation once. A $Z_j$ is one gate controlled on $\mathcal{O}(\log n)$ selector qubits; a reflection is one phase controlled on $n_a$ qubits; a shift by a fixed $\ell$ is $\mathcal{O}(\log|\ell|)$ increments of at most $n_a+1$ qubits, each $n_a+1$ \gls{mcx} gates of up to $n_a$ controls, $\mathcal{O}(n^2)$ gates with a linear-cost decomposition of each. The block-register operations act on $n_f$ qubits, fixed by the problem, and the preparations are rotation trees over $\mathcal{O}(n)$ joint indices. There are $\mathcal{O}(n)$ distinct operations, the $Z_j$ costing $\mathcal{O}(\log n)$ each and the rest $\mathcal{O}(n^2)$ in all, so the total is $\mathcal{O}(n^2)$. A source's rank-one term adds the feature map, of $\mathrm{poly}(n)$ gates, and a rotation tree over its $N$ coefficients, which has $\mathcal{O}(N)$ gates for a general source, as for any general state, and $\mathrm{poly}(n)$ when only a fixed number of the coefficients is non-zero, as for a polynomial source of fixed degree; the further combination that adds it costs $\mathcal{O}(1)$ selector qubits.

\paragraph{Exactness.} \Cref{prop:term-ir-compilation}; every term is a unitary, so no part of the encoding is approximate.

\paragraph{Subnormalisation.} Take first the chain of a regular second-order problem at a fixed $\zeta\ge1$. The recurrences' $\Dt$ terms share the block-register pattern $\operatorname{diag}(1,\zeta,0,0)$, whose Pauli expansion has $\alpha=\tfrac12(|1+\zeta|+|1-\zeta|)=\max(1,\zeta)=\zeta$, its largest entry. With \Cref{lem:affine-tight}, $\Dt$'s affine part is encoded at $2(N-1)$, so these terms contribute $2\zeta(N-1)$ to $\alpha_A$. Every other term is independent of $N$: $\Dt$'s exception at column~1, the $\Bt$ bands weighted by $\zeta$ and $\zeta^2$, the coefficients' multiplication bands weighted by at most $\zeta^2$, and the padding. A source's rank-one term has coefficient $\|\hat r\|\,\|\tau(x_s)\|/|f(x_s)|$, and $\|\hat r\|^2$ grows as $N$ for a fixed smooth source (\Cref{eq:eta-e}), while $\|\tau(x_s)\|^2<2$ (\Cref{prop:tau-norm}). On the other side, $\zeta\Dt$ is a block of $\Adesc$, so $\|\Adesc\|_2\ge2\zeta(N-1)$. Hence $\alpha_A/\|\Adesc\|_2\le1+C/N$ without a source and $1+C'/\sqrt N$ with one, with constants that grow with $\zeta$ and the coefficients. For the mass form of a scalar equation $f'=af$, $\alpha_A\le2(N-1)+\mathcal{O}(1)$ and $\|\Dt-a\Bt\|_2\ge2(N-1)-|a|(2+\sqrt2)$, with the same conclusion. For a chain of order three or more the recurrences' pattern need not be encoded at its largest entry, and no limit is claimed; no benchmark problem has one. On a singular problem $\zeta=N/2$ makes $\zeta\Dt$, $\zeta^2\Bt$ and $\zeta^2a_2$ all of order $N^2$ in one row block, and terms of comparable size add in $\ell^1$ in $\alpha_A$, so the ratio tends to a constant, measured near two (\Cref{sec:descriptor-cost}).

\paragraph{Optimality.} An exact combination whose SELECT indexes $L$ distinct operations needs $\lceil\log_2L\rceil$ selector qubits \cite{chakraborty2025qsvt}. $\Lambda$ has exactly $n+1$ non-zero Walsh coefficients, so a combination of Pauli-$Z$ strings needs $n+1$ terms for it and $\Omega(\log n)$ selector qubits. The descriptor's count is therefore of the least order among exact combinations that expand their diagonals in Pauli strings, and is not claimed against encodings built otherwise.

\paragraph{Norms in \Cref{tab:descriptor-regimes}.} Each row block of $R_{\mathrm{desc}}$ is a finite sum of blocks weighted by powers of $\zeta$: $\zeta^m\Dt$, of norm $2\zeta^m(N-1)$; $\zeta^{m+1}\Bt$, of norm at most $(2+\sqrt2)\zeta^{m+1}$; $\zeta^j\M_{a_j}$, of norm at most $\sqrt2\,\zeta^j\sup|a_j|$ (\Cref{prop:lift}); the padding, of norm one; and the rank-one constraint and source rows, of norm $\mathcal{O}(\sqrt N)$, since $\Brow(x)=\sqrt N\ket0\bra{\tau(x)}$ with $\|\tau(x)\|^2<2$. The norm of $R_{\mathrm{desc}}$ is at most the sum of these and at least the largest of them. At fixed $\zeta$ the largest is $2\zeta^{k-1}(N-1)$, so $\|R_{\mathrm{desc}}\|_2=\Theta(N)$ and $\|\Hdesc\|=\Theta(N^2)$. With $\zeta=N/2$ and $k=2$, $\zeta\Dt$, $\zeta^2\Bt$ and $\zeta^2\M_{a_2}$ are each $\Theta(N^2)$, so $\|\Hdesc\|=\Theta(N^4)$; on two axes with $\zeta_a=N_a$ and $k_a=2$ the same count gives $\Theta(N^4)$; and with $\zeta\propto N^2$ the block $\zeta^2\Bt$, of norm at least $2\zeta^2$, gives $\Theta(N^8)$.

\subsection{Success Probability of One Query}
\label{app:haar-success}

\begin{proposition}[One query on a Haar-random state]
\label{prop:haar-success}
One application of a block encoding of $R$ at $\alpha_R$, post-selected, succeeds on a Haar-random state of dimension $D$ with mean probability $\|R\|_F^2/(D\alpha_R^2)$.
\end{proposition}
\begin{proof}
The success probability on $\psi$ is $\|R\psi\|^2/\alpha_R^2$, and $\mathbb{E}\,\psi\psi^\dagger=I/D$ for Haar-random $\psi$, so its mean is $\operatorname{tr}(R^\dagger R)/(D\alpha_R^2)$.
\end{proof}

In the structured standard form the dense $\G^k$ dominates: for $k=1$, $\|\G\|_F^2\sim N^4/2$ (\Cref{prop:G-structure}) against $\alpha_\G^2\sim N^4$ on $D=N$ amplitudes, so the mean falls as $1/(2N)$. In the descriptor the dominant $\Dt$ has $\|\Dt\|_F^2=\sum_{q<N}(2q)^2\sim4N^3/3$, up to its first entry, against $\alpha^2\sim4N^2$, so the mean stays $\Theta(1)$. The pipeline post-selects no single query, its walk being unitary, so this compares the encodings and is not a cost of the preparation.
